\chapter{Предлагаемый метод} \label{ch3}

В данной главе описан разработанный алгоритм приоритизации тестовых случаев GA-FP+Div, представляющий 
собой генетический алгоритм с улучшенной многокомпонентной фитнесс-функцией.

\section{Фитнесс-функция}\label{ch3:sec1}

Фитнесс-функция оценивает качество предложенной последовательности тестов по трём нормализованным компонентам. 
Использование нескольких компонент позволяет учитывать не только покрытие кода, но и рискованность покрываемых 
методов и структурное разнообразие тестов.

\textbf{Покрытие кода} ($\text{Cov}$) --- доля уникальных методов, покрытых последовательностью тестов, от 
общего числа методов проекта:
$$\text{Cov}(T') = \frac{|\bigcup_{i=1}^{n} M_i|}{|M|}$$

Числитель --- количество уникальных методов, покрытых хотя бы одним тестом из последовательности. 
Знаменатель --- общее число методов в проекте. Результат нормализован к диапазону $[0, 1]$.

\textbf{Вероятность дефектов} ($\text{FP}$) --- средняя вероятность дефектов первых $k$ тестов последовательности:
$$\text{FP}(T') = \frac{1}{k}\sum_{i=1}^{k} r_{\pi(i)}$$

где $r_{\pi(i)}$ --- вероятность дефектов $i$-го теста в последовательности, оцениваемая на основе анализа 
влияния изменений CIA (Change Impact Analysis). Компонента максимизирует вероятность того, что наиболее 
рискованные тесты выполняются первыми.

\textbf{Структурное разнообразие} ($\text{Div}$) --- среднее расстояние Жаккара между соседними тестами в 
первых $k$ позициях:
$$\text{Div}(T') = \frac{1}{k-1}\sum_{i=1}^{k-1} d_J(M_{\pi(i)}, M_{\pi(i+1)})$$

где расстояние Жаккара определяется как:
$$d_J(A, B) = 1 - \frac{|A \cap B|}{|A \cup B|}$$

Значение $d_J = 0$ означает что тесты покрывают одинаковые методы, $d_J = 1$ --- покрытия не пересекаются. 
Компонента стимулирует алгоритм чередовать тесты с различным покрытием, снижая дублирование в начале последовательности.

Все три компоненты нормализованы к диапазону $[0, 1]$. Итоговый фитнесс вычисляется как взвешенная сумма:
$$f(T') = w_c \cdot \text{Cov}(T') + w_r \cdot \text{FP}(T') + w_d \cdot \text{Div}(T')$$

где $w_c + w_r + w_d = 1$, $w_c, w_r, w_d \geq 0$. В настоящей 
работе приняты значения $w_c = 0{,}2$, $w_r = 0{,}6$, $w_d = 0{,}2$, 
$k = 10$.

\section{Операторы генетического алгоритма} \label{ch3:sec2}

Генетический алгоритм оперирует популяцией из $N$ особей, каждая из которых представляет собой перестановку 
тестового набора $T$. На каждом поколении популяция обновляется с помощью операторов селекции, кроссовера и мутации.

\textbf{Инициализация.} Начальная популяция из $N$ особей формируется случайным перемешиванием тестового набора $T$.
Использование фиксированного зерна генератора случайных чисел обеспечивает воспроизводимость результатов.

\textbf{Селекция.} Используется турнирная селекция с размером турнира $s = 3$: случайно выбираются $s$ особей, 
из которых в следующее поколение передаётся особь с наибольшим значением фитнесса. Турнирная селекция обеспечивает 
селективное давление без полного исключения слабых особей.

\textbf{Кроссовер.} Применяется оператор OX (Order Crossover), сохраняющий валидность перестановки. 
Алгоритм OX работает следующим 
образом:
\begin{enumerate}
    \item Случайно выбирается сегмент $[\alpha, \beta]$ из первого 
    родителя $P_1$ и копируется в потомка $C$ на те же позиции.
    \item Оставшиеся позиции $C$ заполняются элементами из второго 
    родителя $P_2$ в порядке их появления, пропуская элементы уже 
    присутствующие в сегменте.
\end{enumerate}

Например, для $P_1 = [t_3, t_1, t_5, t_2, t_4]$ и 
$P_2 = [t_2, t_4, t_1, t_5, t_3]$ при сегменте $[1, 3]$:
$$C = [t_2, \underbrace{t_1, t_5, t_2}_{\text{сегмент из } P_1}, t_3]
\rightarrow [t_4, t_1, t_5, t_2, t_3]$$

OX гарантирует что потомок является корректной перестановкой без дубликатов, что критично для задачи TCP.

\textbf{Мутация.} Выполняется $\lfloor\sqrt{n}\rfloor$ случайных перестановок пар элементов (swap), 
где $n$ --- размер тестового набора. Адаптивная интенсивность мутации позволяет сохранять разнообразие популяции 
при различных размерах тестовых наборов: для $n = 100$ выполняется 10 перестановок, для $n = 1000$ --- 31.

\textbf{Элитизм.} Лучшая особь текущего поколения переносится в следующее поколение без изменений. 
Это гарантирует монотонное неубывание лучшего найденного решения на протяжении всей эволюции.

\textbf{Параметры алгоритма.} Вероятность кроссовера $p_c = 0{,}8$ и вероятность мутации 
$p_m = 0{,}2$ выбраны в соответствии с рекомендациями~\cite{paygude2020}. Размер популяции $N = 100$. 
Число поколений задаётся адаптивно: 
$G = \max(50,\ \lfloor\log_2(n)\rfloor \cdot 10)$.

\section{Псевдокод алгоритма} \label{ch3:sec3}
Алгоритм работает следующим образом. На этапе инициализации формируется популяция $P$ из $N$ 
случайных перестановок тестового набора $T$, из которых выбирается лучшая особь $T^*$.

В каждом из $G$ поколений формируется новая популяция $P'$. Первой в неё помещается текущая лучшая 
особь $T^*$ без изменений (элитизм). Затем в цикле до заполнения популяции выполняются следующие шаги: 
турнирной селекцией выбираются два родителя $P_1$ и $P_2$; с вероятностью $p_c$ к ним применяется 
OX-кроссовер, порождающий двух потомков $C_1$ и $C_2$; с вероятностью $p_m$ каждый потомок подвергается 
мутации. Сформированные потомки добавляются в $P'$.

После обновления популяции проверяется улучшение: если лучшая особь нового поколения превосходит 
$T^*$ по значению фитнесса, $T^*$ обновляется. По завершении всех поколений алгоритм возвращает 
$T^*$ --- наилучшую найденную перестановку тестового набора.

\begin{lstlisting}[caption={Псевдокод алгоритма GA-FP+Div}, 
                   label={lst:ga}]
Вход: T, N, G, pc, pm, wc, wr, wd, k
Выход: T* -- оптимальная перестановка

P <- InitPopulation(T, N)
T* <- argmax f(P)

для g от 1 до G:
    P' <- {T*}                        // элитизм
    пока |P'| < N:
        P1 <- Tournament(P, 3)
        P2 <- Tournament(P, 3)
        если rand() < pc:
            C1, C2 <- OXCrossover(P1, P2)
        иначе:
            C1 <- P1; C2 <- P2
        если rand() < pm:
            Mutate(C1, floor(sqrt(n)))
        если rand() < pm:
            Mutate(C2, floor(sqrt(n)))
        P' <- P' + {C1, C2}
    P <- P'
    если f(best(P)) > f(T*):
        T* <- best(P)

вернуть T*
\end{lstlisting}